\begin{frame}{Колоризація}
        \begin{columns}[T]
            \begin{column}{0.65\textwidth}
                % \textbf{Колоризація зображень}

                % Це процес відновлення зображення або додавання кольору до чорно-білого зображення.
                % % \begin{itemize}
                % %     \item Це процес відновлення зображення або додавання кольору до чорно-білого зображення.
                % %     % \item \textbf{Математично:} Прогнозування хроматичних каналів ($a$ та $b$) на основі каналу яскравості ($L$) у колірному просторі \textbf{CIELAB}.
                % % \end{itemize}
                
                % \vspace{1em}
                % \textbf{Користувацькі підказки}

                % Це кольорові точки, які допомагають вирішити неоднозначність, наприклад, колір одягу зробити зеленим.
                % % \begin{itemize}
                % %     \item Це кольорові точки, які допомагають вирішити неоднозначність, наприклад, колір одягу зробити зеленим.
                % %     % \item \textbf{Формат:} Задаються як м'які Гауссівські плями, а не жорсткі пікселі. Це дозволяє моделі плавно поширювати колір без різких меж та артефактів.
                % % \end{itemize}
                \begin{block}{Колоризація зображень}
                    \begin{itemize}
                        \item Це процес відновлення зображення або додавання кольору до чорно-білого зображення.
                        % \item \textbf{Математично:} Прогнозування хроматичних каналів ($a$ та $b$) на основі каналу яскравості ($L$) у колірному просторі \textbf{CIELAB}.
                    \end{itemize}
                \end{block}
                
                % \vspace{1em}
                \begin{exampleblock}{Користувацькі підказки}
                    \begin{itemize}
                        \item Це кольорові точки, які допомагають вирішити неоднозначність, наприклад, колір одягу зробити зеленим.
                        % \item \textbf{Формат:} Задаються як м'які Гауссівські плями, а не жорсткі пікселі. Це дозволяє моделі плавно поширювати колір без різких меж та артефактів.
                    \end{itemize}
                \end{exampleblock}
            \end{column}

            \begin{column}{0.35\textwidth}
                \centering
                % \vspace{1.0em}
                \vspace{0.8em}
                \begin{tikzpicture}[
                    arrow/.style={->, ultra thick, >=stealth, color=white}
                ]

                    \node[inner sep=0pt, draw, thick] (imgA) at (0, 1.25) {
                        \includegraphics[width=3.5cm, height=2cm]{CharlieChaplin_gray_hint.png}
                    };
                    \node[inner sep=0pt, draw, thick] (imgB) at (0, -1.25) {
                        \includegraphics[width=3.5cm, height=2cm]{CharlieChaplin_color.png}
                    };
                    \draw[arrow] (imgA.south) -- (imgB.north);

                \end{tikzpicture}
            \end{column}
        \end{columns}
    \end{frame}